% Generated draft: requires mathematical and source review.
\documentclass[11pt]{article}
\usepackage[margin=25mm]{geometry}
\usepackage{fontspec}
\setmainfont{Latin Modern Roman}
\usepackage{amsmath,amssymb,mathtools,unicode-math}
\setmathfont{Latin Modern Math}
\usepackage{enumitem,needspace,xurl,hyperref}
\setlength{\parindent}{0pt}
\setlength{\parskip}{5pt}
\setlength{\emergencystretch}{3em}
\newcommand{\N}{\mathbb{N}}
\newcommand{\Z}{\mathbb{Z}}
\newcommand{\Q}{\mathbb{Q}}
\newcommand{\R}{\mathbb{R}}
\newcommand{\C}{\mathbb{C}}
\newcommand{\F}{\mathbb{F}}
\newcommand{\A}{\mathbb{A}}
\newcommand{\PP}{\mathbb{P}}
\newcommand{\E}{\mathbb{E}}
\newcommand{\eps}{\varepsilon}
\newcommand{\dd}{\,\mathrm{d}}
\begin{document}
\section*{UnsolvedMath ID: 30001905}
\textbf{Source problem code:} catalogue entry
\subsection{Definitions and mathematical statement}
Let \(S\) denote the sphere spectrum, let \(MU\) denote the complex bordism spectrum, and let \(S\to MU\) be its unit map. Write \(K(A)\) for the algebraic \(K\)-theory spectrum of a ring spectrum \(A\); the supplied record does not specify the version of algebraic \(K\)-theory.

The source asks whether \(K(S)\) can be recovered from \(K(MU)\) by descent along \(S\to MU\), which it describes as a Hopf--Galois extension. The original wording also asks whether \(K(MU)\) is simpler to describe than \(K(S)\).

One possible precise interpretation of the descent question is whether the canonical comparison
\[
K(S)\longrightarrow
\operatorname*{Tot}_{[q]\in\Delta}
K\!\left(MU^{\wedge_S(q+1)}\right)
\]
is an equivalence of spectra. Here \(\wedge_S\) denotes the derived smash product over \(S\), \(\Delta\) is the simplex category, and \(\operatorname{Tot}\) denotes the homotopy limit of the cosimplicial spectrum obtained by applying \(K\) to the Amitsur construction of \(S\to MU\). Its coface maps insert the unit and its codegeneracy maps multiply adjacent factors.

This displayed comparison is an explicitly unconfirmed interpretation, not a formulation specified by the supplied source record. The record does not identify the Hopf object, coaction, descent construction, or any completion or localization intended in the recovery question. Consequently, it does not determine a unique precise descent assertion. The additional phrase ``simpler to describe'' has no supplied mathematical criterion.
\subsection{Short English statement}
The problem asks whether algebraic \(K\)-theory of the sphere spectrum can be reconstructed from algebraic \(K\)-theory of complex bordism using the descent structure associated with \(S\to MU\). The exact descent construction is unspecified; the displayed Amitsur comparison is only a possible interpretation. The original source also asks, qualitatively, whether \(K(MU)\) is simpler to describe than \(K(S)\).
\subsection{Sources}
\begin{itemize}
\item Supplied local UnsolvedMath registry record, permanent ID \(30001905\), source problem code \texttt{catalogue entry}, title ``K-Theoretic Descent from MU to the Sphere''; registry collection \texttt{data/unsolvedmath/problems.json}. The supplied record belongs to the Oberwolfach Reports Open Problems collection and records proposed year \(2011\). No individual proposer is supplied.
\item Supplied report citation: \emph{Homotopy Theory (2012)}, Oberwolfach Reports, pp. 2623--2667, DOI \texttt{10.4171/owr/2011/46}. Source URL and literature source labeled ``original OWR report'': \url{https://doi.org/10.4171/owr/2011/46}. The background describes Oberwolfach Reports as documenting mathematical workshops and their open-problem discussions.
\item Supplied literature assessment, checked on \texttt{2026-08-21}: the record labels the problem open and states that the descent question remains a research problem, with the report providing motivation but no complete descent or reconstruction theorem. This is supplied dated curation derived from the OWR record, not an independently verified literature assessment. The source itself requests independent verification of time-sensitive status and cited claims.
\item The registry labels its cleaned statement \texttt{corrected\_verified} and reports a manual semantic, typographic, and LaTeX review. These are supplied editorial claims, not independent verification of the descent formulation or mathematical status.
\end{itemize}
\end{document}
